\begin{theorem}[Preparation of the PE special closure]
Let \(G\) be a finite loopless multigraph with its exact edge type, let
\(e_0\) be the distinguished edge, and let the palette have \(k\) colors.
Let
\[
 \mathcal E:\noderef{ExtendedTashkinovTree}(G,k,e_0,\rho,T,\varphi)
\]
be a maximum ETT, put \(n=\mathcal E.\mathit{rounds}>0\), and put
\[
 a_n=\operatorname{Fin.last}(n),\qquad
 T_n=\mathcal E.\mathit{series.data.tree}(a_n.\mathit{castSucc}).
\]
Assume that the extension kind at \(a_n\) is PE.  Let
\(B\) be a generating-color \(\noderef{HierarchyBase}\) for
\((\mathcal E,\varphi)\), and assume
\[
 \noderef{MaximumDegree}(G)+1\leq k,
 \qquad
 2n+11\leq |\noderef{TreeMissingColors}(\varphi,T_n)|.          \tag{1}
\]
Suppose also that a
\(\noderef{TechnicalInductionContext}(G,k,e_0,n)\) is given.

Then
\[
 \noderef{TashkinovAugmentationSequence}
       (\varphi,T_n,B.\mathit{core}),
 \qquad |T_n|<|B.\mathit{core}|.                                \tag{2}
\]
Moreover, one of the following two alternatives holds.

First, \(V(T)\) is
\(\noderef{ColoredVertexSetElementary}\) under \(\varphi\).

Second, there are an ordered tree-sequence \(U\) and a certificate
\[
 C_U:\noderef{CorrespondingExtendedTashkinovTree}
       (\mathcal E,\varphi,U)                                   \tag{3}
\]
relative to the fixed original ETT \(\mathcal E\), such that the replacement
ETT \(\mathcal E_U=C_U.\mathit{replacementETT}\) has maximum property and
\[
 \noderef{TashkinovAugmentationSequence}
       (\varphi,B.\mathit{core},U),\qquad
 B.\mathit{core}\preccurlyeq U,
 \qquad |B.\mathit{core}|<|U|.                                  \tag{4}
\]
Here \(\preccurlyeq\) is the exact ordered
\(\noderef{TreeSequenceInitialSegment}\) relation.

There are a hierarchy base \(B_U\) for \((\mathcal E_U,\varphi)\), a
\(\noderef{TashkinovHierarchy}\) \(K_U\) based on \(B_U\), and a
\(\noderef{GoodHierarchy}\) witness \(\mathcal G_U\) such that
\[
 \begin{split}
 B_U.\mathit{core}&=B.\mathit{core},\\
 B_U.\mathit{closure}&=B.\mathit{closure},\\
 B_U.\mathit{supportPosition.val}&=B.\mathit{supportPosition.val},
 \end{split}                                                     \tag{5}
\]
and
\[
 K_U.\mathit{dividerCount}=0,\qquad K_U(0)=T_n,
 \qquad K_U(\operatorname{last}(K_U.\mathit{dividerCount}+1))=U. \tag{6}
\]
Write \(D_0=\noderef{HierarchyStoredColorsAt}(K_U,0)\), write
\(\Gamma^0_\eta=\mathcal G_U.\mathit{gamma}(0,\eta)\), and put
\(\Gamma^0=\bigcup_{\eta\in D_0}\Gamma^0_\eta\).  For every
\(\eta\in D_0\),
\[
 |\Gamma^0_\eta|=2,
 \qquad
 \Gamma^0_\eta\subseteq
 \overline\varphi(T_n)-
   \noderef{HierarchySegmentColors}(K_U,0,\eta),                 \tag{7}
\]
and the sets \(\Gamma^0_\eta\) for distinct members of \(D_0\) are
pairwise disjoint.  In addition,
\[
 \begin{split}
 &B.\mathit{core}\text{ is not }(\Gamma^0)^-\text{-closed},     \tag{8}\\
 &\left|\bigl(\overline\varphi(T_n)\cap
       \overline\varphi(B.\mathit{closure})\bigr)-\Gamma^0\right|
       \geq4.                                                    \tag{9}
 \end{split}
\]
Finally, with the terminal stored-color set kept distinct from \(D_0\),
\[
 U\text{ is }
 \left(\bigcup_{\eta\in
      \noderef{HierarchyTerminalStoredColors}(K_U)}
      \Gamma^0_\eta\right)^-\text{-closed under }\varphi.      \tag{10}
\]
All tree equalities and initial-segment relations above concern the full
ordered edge-valued tree-sequences.
\end{theorem}
\begin{proof}
We keep the original ETT fixed throughout the construction.  In particular,
every occurrence of ``corresponding'' in the conclusion is relative to
\(\mathcal E\), rather than merely to an ETT obtained at an intermediate
replacement step.

First assemble the lower-rung hypothesis expected by the existing technical
results.  Given an arbitrary positive-rung maximum ETT \(\mathcal L\) with
\(\mathcal L.\mathit{rounds}<n\), apply, in order, the six fields
\(\mathit{lowerOne},\ldots,\mathit{lowerSix}\) of the assumed
\noderef{TechnicalInductionContext}.  Their conclusions have exactly the
same independently quantified input coloring, displayed tree, generating
coloring, and ETT.  Conjoin them without changing those parameters.  Denote
the resulting six-assertion conjunction by \(H_{<n}(\mathcal L)\).  Thus its
first component is lower-rung elementarity, its second is interchangeability
on the exact final series tree (including isolated path components), its
third contains both stable-before-first-appearance and raw stored-set
disjointness, its fourth is the PE two-color-component assertion, its fifth
is the defective-vertex order assertion, and its sixth is the lower-rung
mod-coloring assertion.  The positivity condition in the predecessor
stability clause is retained verbatim; no zero-rung ETT is introduced.

We next record the terminal coloring facts.  The positivity field of
\noderef{ExtendedTashkinovTree} makes \(a_n\) a positive series index.  The
certificate at that index supplied by \noderef{TashkinovSeries}, together
with the assumed PE kind, selects the PE constructor of
\noderef{TashkinovSeriesStep}.  Its current-coloring, current-properness, and
current-one-uncolored clauses, combined with the generating-coloring equality
in the ETT, show that \(\varphi\) is proper and leaves exactly \(e_0\)
uncolored.  The same constructor also says that the original final series
tree
\[
 F=\mathcal E.\mathit{series.data.tree}
       (\operatorname{Fin.last}(n+1))                            \tag{11}
\]
is a \(\varphi\)-\noderef{TashkinovClosure} of \(T_n\).  Since stability of
a coloring with itself is the reflexive clause of \noderef{sc2}, \(\varphi\)
is stable on \(T_n\), relative to the computed terminal
\noderef{TashkinovStoredColors}, with respect to itself.

Apply \noderef{Td} to \(\mathcal E\), its maximum property, its positive rung,
the locally assembled conjunction \(H_{<n}\), and this proper,
exactly-one-uncolored, reflexively stable generating coloring.  We obtain
\[
                         M_0:\noderef{ModColoring}(\mathcal E,\varphi).
                                                                    \tag{12}
\]
Independently, \noderef{ParallelCoreAugmentationStrict}, applied to the PE
kind and the base \(B\), gives exactly (2).  Call its augmentation sequence
\(P:T_n\longrightarrow B.\mathit{core}\).

We now establish the three PE facts that will later be supplied to Algorithm
5.5.  Apply \noderef{ParallelClosureUnionElementary} with
\(H_{<n}\), the maximum property, the PE kind, the coloring
\(\sigma=\varphi\), reflexive stability, and the given base.  It yields
\[
 V(B.\mathit{core})\text{ is elementary under }\varphi.          \tag{13}
\]
The input required for common-missing-color closedness is the elementarity of
the exact predecessor ladder tree under the pre-step coloring.  We derive it
without postulating a zero-rung ETT.

If \(n=1\), the first conclusion of \noderef{TechnicalTheoremBase}, applied
to the original certified series, states that its exact tree at data index
one is elementary under the coloring at index zero.  After substituting
\(n=1\), those are precisely \(T_n\) and its predecessor coloring.  If
\(1<n\), apply \noderef{CertifiedETTPredecessor} to \(\mathcal E\).  It
produces an actual positive-rung maximum predecessor \(\mathcal L\), with
displayed tree exactly \(T_n\), generating coloring exactly the predecessor
coloring, and rung number \(n-1<n\).  The field
\(\mathit{lowerOne}\) of the context, applied to this certified
\(\mathcal L\), proves the same elementarity assertion.  These two cases
therefore supply the precise premise of \noderef{ParallelCommonMissingClosed},
which gives
\[
 \noderef{TreeColorSetClosed}\left(
   G,\varphi,B.\mathit{core},
   \overline\varphi(T_n)\cap
     \overline\varphi(B.\mathit{closure})\right).                \tag{14}
\]
Finally, apply \noderef{ParallelBaseFourCommonMissingColors} to the proper
exactly-one-uncolored coloring, the degree bound, and (13).  We obtain
\[
 4\leq\left|\overline\varphi(T_n)\cap
              \overline\varphi(B.\mathit{closure})\right|.       \tag{15}
\]

Choose, using \noderef{TashkinovClosureExists}, a
\(\varphi\)-Tashkinov closure \(Z\) of \(B.\mathit{core}\).  Its closure
certificate contains an augmentation sequence
\(Q:B.\mathit{core}\longrightarrow Z\).  By the reflexive-transitive-chain
definition in \noderef{TashkinovAugmentationSequence}, concatenate \(P\) and
\(Q\) to obtain a sequence \(T_n\longrightarrow Z\).  Because (2) is strict,
\(|T_n|<|Z|\); and
\noderef{TashkinovAugmentationSequenceInitialSegment} makes \(T_n\) an exact
ordered initial segment of \(Z\).  The universal PE clause of the
mod-coloring (12) therefore gives a certificate
\[
 C_0:\noderef{CorrespondingExtendedTashkinovTree}
        (\mathcal E,\varphi,Z).                                  \tag{16}
\]
Let \(\mathcal E_0=C_0.\mathit{replacementETT}\).  The theorem
\noderef{MaximumPropertyTransport}, applied to (12), (16), and the original
maximum property, gives maximum property for \(\mathcal E_0\).
The theorem \noderef{HierarchyBaseCorrespondingTransport} gives a base
\(B_0\) for \((\mathcal E_0,\varphi)\) satisfying
\[
 B_0.\mathit{core}=B.\mathit{core},\quad
 B_0.\mathit{closure}=B.\mathit{closure},\quad
 B_0.\mathit{supportPosition.val}=B.\mathit{supportPosition.val}. \tag{17}
\]

It is essential here to identify the complete tree \(Z\), not merely its
vertex set.  Let
\[
 F_0=\mathcal E_0.\mathit{series.data.tree}
       (\operatorname{Fin.last}(\mathcal E_0.\mathit{rounds}+1))  \tag{18}
\]
be the exact final series output of the replacement ETT.  The rung,
ladder-tree, and extension-kind agreement fields in (16) identify the
replacement terminal ladder tree with \(T_n\) and its terminal kind with PE.
Consequently its terminal certified \noderef{TashkinovSeriesStep} is the PE
constructor, and that constructor makes \(F_0\) a
\(\varphi\)-Tashkinov closure of the same exact tree \(T_n\).

On the other hand, the concatenated route \(T_n\longrightarrow Z\), together
with the closedness part of the chosen closure of
\(B.\mathit{core}\), makes \(Z\) itself a
\(\varphi\)-Tashkinov closure of \(T_n\).  The `belowFinal' field of
\(\mathcal E_0\) says that its displayed tree \(Z\) is an exact ordered
initial segment of \(F_0\).  Apply
\noderef{TashkinovClosureVertexSetUnique} to the two closures from the exact
same source \(T_n\); it gives
\(V(Z)=V(F_0)\).  Then
\noderef{TreeSequenceInitialSegmentVertexSetEq} upgrades this equality and
\(Z\preccurlyeq F_0\) to the equality of complete ordered edge-valued trees
\[
                              Z=F_0.                              \tag{19}
\]
Thus, and only now, rewriting the original closure
\(B.\mathit{core}\longrightarrow Z\) by (17) and (19) supplies the closure
\[
              B_0.\mathit{core}\longrightarrow F_0              \tag{20}
\]
required for the replacement presentation.

We must derive a second mod-coloring rather than reuse (12).  For every
positive-rung maximum ETT \(\mathcal L\) with
\(\mathcal L.\mathit{rounds}<\mathcal E_0.\mathit{rounds}\), the rung
agreement in (16) changes this inequality into
\(\mathcal L.\mathit{rounds}<n\).  Apply the six lower fields of the same
\noderef{TechnicalInductionContext} and conjoin them exactly as before;
call the result \(H^0_{<n}(\mathcal L)\).  The type of (16) says that the
generating coloring of \(\mathcal E_0\) is again \(\varphi\).  Its terminal
properness and exactly-one-uncoloredness follow from the replacement
properness fields of (16), while reflexive stability again follows from
\noderef{sc2}.  Apply \noderef{Td} to \(\mathcal E_0\), its transported
maximum property, its positive rung transported through the rung agreement,
and \(H^0_{<n}\).  This gives
\[
                         M_1:\noderef{ModColoring}(\mathcal E_0,\varphi).
                                                                    \tag{21}
\]

We now split according as the chosen core closure grew.  Suppose first that
\(Z=B.\mathit{core}\).  Then the closedness of \(Z\) and the route \(P\)
make \(B.\mathit{core}\) a closure of \(T_n\).  The original terminal PE
certificate made \(F\) in (11) another closure of \(T_n\).  By
\noderef{TashkinovClosureVertexSetUnique},
\[
                        V(B.\mathit{core})=V(F).                  \tag{22}
\]
The `belowFinal' field of the original
\noderef{ExtendedTashkinovTree} gives \(T\preccurlyeq F\); expanding
\noderef{TreeSequenceInitialSegment} and
\noderef{TreeSequenceVertexSet} gives \(V(T)\subseteq V(F)\).  After (22),
this is \(V(T)\subseteq V(B.\mathit{core})\).  Elementarity in (13) is
hereditary under restriction: two distinct vertices of \(V(T)\) are two
distinct vertices of the elementary superset, so their missing-color sets
are disjoint by the definition of \noderef{ColoredVertexSetElementary}.
Hence \(V(T)\) is elementary, which is the first alternative.

It remains to treat \(Z\ne B.\mathit{core}\).  The initial-segment relation
implicit in \(Q\), supplied explicitly by
\noderef{TashkinovAugmentationSequenceInitialSegment}, first gives
\(|B.\mathit{core}|\leq|Z|\).  Equality of these lengths would make its
vertex and edge equalities cover every index of both trees; extensionality of
the data fields in \noderef{TreeSequence}, followed by proof irrelevance for
the structural fields, would then give \(Z=B.\mathit{core}\).  This contradicts
the present case.  Hence
\(|B.\mathit{core}|<|Z|\).  Apply
\noderef{TashkinovAugmentationSequenceStrictBoundaryWitness} to \(Q\).
It supplies a color \(\alpha\) missing on \(B.\mathit{core}\) and a concrete
\(\alpha\)-colored boundary edge identity.  By the meanings of
\noderef{BoundaryColorEdges} and \noderef{TreeColorSetClosed}, this says
\[
 \alpha\in\overline\varphi(B.\mathit{core}),\qquad
 \neg\,\noderef{TreeColorSetClosed}
       (G,\varphi,B.\mathit{core},\{\alpha\}).                    \tag{23}
\]

We verify the inputs of \noderef{AlgorithmFiveFiveBaseStage} for
\(\mathcal E_0,B_0\).  The replacement terminal kind is PE by (16).  Its
terminal ladder tree is exactly \(T_n\); (17) transports the strict
inequality in (2) to \(|T_n|<|B_0.\mathit{core}|\).  Equation (20) is its
required core-to-final closure.  Estimate (1) transports across the exact
ladder-tree equality in (16).  Assertion (23), assertion (14), and assertion
(15), rewritten using (17) and the same ladder-tree equality, are exactly the
three PE facts required by that theorem.  Finally use the maximum property of
\(\mathcal E_0\) and its independently derived mod-coloring (21).

The output of \noderef{AlgorithmFiveFiveBaseStage} consists of a seed \(J\),
a terminal tree \(U\), and a certificate
\[
 C_1:\noderef{CorrespondingExtendedTashkinovTree}
       (\mathcal E_0,\varphi,U).                                 \tag{24}
\]
Because the terminal kind is PE, its exact seed alternative says
\(J=B_0.\mathit{core}=B.\mathit{core}\).  Hence its augmentation route,
initial-segment clause, and strict PE growth give exactly (4).  It also gives
maximum property for \(C_1.\mathit{replacementETT}\), a base \(B_1\), a
zero-divider hierarchy \(K_1\), and a good witness \(\mathcal G_1\).  Its
displayed conclusions, after (17) and the ladder equality from (16), are
precisely (5)--(10), except that (24) is at this moment relative to the
intermediate ETT \(\mathcal E_0\).

Apply \noderef{CorrespondingExtendedTashkinovTreeTransTrees} immediately to
(16) and (24).  It gives a composite certificate \(C_U\) as in (3), relative
to the fixed original \(\mathcal E\), equality of the replacement inputs,
and heterogeneous equalities between the composite replacement series and
ETT and those supplied by \(C_1\).  Transport maximum property along the ETT
equality.  Transport \(B_1\) along the same equality to obtain \(B_U\), then
transport \(K_1\) along the transported base equality to obtain \(K_U\), and
finally transport \(\mathcal G_1\) along the hierarchy equality to obtain
\(\mathcal G_U\).  These are dependent transports, not an appeal to equality
of vertex sets.

The transported core, closure, and numerical support-position equations are
still (5).  The zero-divider equation, the two full stage equalities, every
stored set, every Gamma pair, and every hierarchy segment color set are
preserved by these dependent equalities.  In the first stage equality, the
ladder agreement in (16) rewrites the intermediate terminal ladder tree to
the original exact \(T_n\), giving (6).  Consequently the pair-size and
avoidance assertions become (7), pairwise disjointness is unchanged, the two
PE conclusions become (8)--(9), and the transported terminal
minus-closedness assertion is exactly (10).  The certificate (3) now has the
required fixed-original provenance, completing the strict branch and the
proof.
\end{proof}
